\subsection{Gluing small triangles along an acyclic corner or edge}

\begin{lemma}[Proved: boundedness and products through a low-homology intersection]
\label{CP2:A-low-homology-gluing}
Suppose full subcategories $E,F$ of a finite category $C$
cover its nerve, and their intersection $I$ satisfies
$H_j(BI;\mathbb Q)=0$ for every $j\geq2$.
Then restriction gives an isomorphism
\[
 H^n(BC;\mathbb Q)\longrightarrow
 H^n(BE;\mathbb Q)\oplus H^n(BF;\mathbb Q)
 \quad(n\geq3)
\]
and a surjection in degree two. Consequently bounded
rational homology of $BC$ passes to both pieces, and
each piece's second Betti number is at most that of $BC$.
If both pieces have zero rational degree-two products,
then so does $BC$.
\end{lemma}
\begin{proof}
The realizations are CW subcomplexes and form a pushout
along their intersection. Apply the rational cohomology
Mayer--Vietoris sequence. For $n\geq3$ both adjacent
groups $H^{n-1}(BI;\mathbb Q)$ and $H^n(BI;\mathbb Q)$
are zero, giving the isomorphism. The group after
$H^2(BE)\oplus H^2(BF)$ is $H^2(BI)=0$, giving the
surjection. Over a field, the corresponding homology
and cohomology vanishing statements agree. These facts
give boundedness and the asserted dimension inequalities.
Finally restriction is multiplicative and is injective
in degree four. A degree-two product whose restrictions
to both pieces vanish must therefore vanish on $BC$.
\end{proof}

In the next two statements a \emph{permitted triangle}
means a full-linear three-object category with profile
$(a,b,c;x,y,z)$ satisfying at least one of these
cardinality conditions:
\[
\begin{gathered}
 z=1,\quad \max(a,b,c)\leq4;\\
 a,c\leq2,\quad b\leq4;\\
 a,b\leq2,\ c\leq6,\quad\hbox{or}\quad b,c\leq2,\ a\leq6;\\
 a\leq2,\ x=1,\ c\leq11,\quad\hbox{or}\quad
 c\leq2,\ y=1,\ a\leq11;\\
 a\leq4,\ x,z\leq2,\quad\hbox{or}\quad
 c\leq4,\ y,z\leq2.
\end{gathered}
\]
Every such triangle with bounded rational homology and
second Betti number at most one has zero degree-two
products. For the five lines, respectively, apply
Theorem~\ref{CP2:A-baseline-singleton-long},
Lemma~\ref{CP2:A-high-linear-Tor},
Corollary~\ref{CP2:A-baseline-semisimple-pair-six},
Corollary~\ref{CP2:A-eleven-projective-singleton}, and
Lemma~\ref{CP2:A-high-small-binary}. The last use takes
the indicated endpoint as one block and the other two
objects as the arbitrary opposite block. Its second-rank
premise is the stated Betti bound; polarization over
$\mathbb Q$ converts zero squares to zero pairwise products.
These are the same underlying triangle statements used
in Corollary~\ref{CP2:A-pendant-bounded-triangle-gates}.

\begin{theorem}[Proved: a pendant edge attached to a permitted triangle]
\label{CP2:A-pendant-edge-permitted-triangle}
Let $C$ have four objects and acyclic support. Suppose
one object $p$ is strictly adjacent only to $q$, the
corner at $q$ has order at most four, and the retained
three objects form a permitted triangle. Let $s$ be the
size of the nonempty hom-set between $p$ and $q$.
If $s\leq2$ or the corner at $p$ also has order at most
four, bounded rational homology and second Betti number
at most one force every rational degree-two product of
$BC$ to be zero.
\end{theorem}
\begin{proof}
Let $E=C\setminus\{p\}$ and $F=C|_{\{p,q\}}$. The absence
of any other strict arrow incident to $p$ ensures that
every nerve simplex lies in $E$ or $F$: a strict chain
containing $p$ and a third object would compose to such
an additional arrow. The intersection is the corner at
$q$, whose classifying space is rationally acyclic by
the small-monoid lemma. Therefore
Lemma~\ref{CP2:A-low-homology-gluing} applies. In fact
here restriction is an isomorphism also in degree two,
since the intersection has zero first rational cohomology.
The permitted triangle $E$ inherits bounded homology
and second Betti number at most one, so its degree-two
products vanish by the five triangle statements above.

For the attached two-object piece $F$, a singleton stalk
is a categorical join and has zero positive products.
A two-element stalk has zero degree-two products by
Lemma~\ref{CP2:A-high-two-cross}, with no corner assumption.
In the remaining case both corner orders are at most
four, so both corner second rational cohomology groups
are zero. The two-block lemma again kills every
degree-two product on $BF$, regardless of its stalk size.
The degree-four restriction injection in the gluing
lemma now kills every such product on $BC$.
\end{proof}

\begin{theorem}[Proved: two permitted triangles meeting in a small edge]
\label{CP2:A-two-permitted-triangles}
Let $C$ have four objects and acyclic support. Suppose
two objects $p,r$ are incomparable, both full subcategories
$E=C\setminus\{p\}$ and $F=C\setminus\{r\}$ are
permitted triangles, and their two-object intersection
$I$ has a nonempty strict hom-set. Assume either that
both corner orders in $I$ are at most two, or that its
strict hom-set is singleton and at least one of its
corner orders is at most four. Bounded rational homology
and second Betti number at most one then force every
rational degree-two product on $BC$ to be zero.
\end{theorem}
\begin{proof}
A nerve simplex cannot contain both $p$ and $r$, since
composing its strict arrows would make these objects
comparable. Thus $BE$ and $BF$ cover $BC$ with intersection
$BI$. If both intersection corners have order at most
two, their rational algebras are semisimple. The relative
bar complex of the two-object edge has strict height
one, hence $H_j(BI;\mathbb Q)=0$ for $j\geq2$.
In the other case $BI$ is the join of its corner spaces;
at least one is rationally acyclic, so their join is
rationally acyclic as well.

The gluing lemma passes bounded homology and the
second-rank bound to each triangle. The five triangle
statements give zero degree-two products on both pieces.
Injectivity of degree-four restriction gives the result.
\end{proof}

\begin{corollary}[Proved: actual deletion of a contractible-neighbor pendant]
\label{CP2:A-contractible-pendant-deletion}
If a singleton pendant $p$ has neighbor $q$ with contractible
corner classifying space, the inclusion
$B(C\setminus\{p\})\to BC$ is an actual homotopy equivalence.
In particular an eighteen-arrow projective-plane model
cannot have such a pendant when the neighbor's minimal
group is trivial and its minimal ideal has one row or
one column.
\end{corollary}
\begin{proof}
The same nerve union has a contractible intersection and
contractible attached join, so the inclusion of the
retained piece is a homotopy equivalence. A monoid with
one minimal row or one minimal column is homotopy
equivalent to its minimal group, by
Lemma~\ref{CP2:A-extreme-minimal-group-reduction} applied
to its one-object category or its opposite. A trivial
minimal group therefore makes the corner contractible.
Deleting $p$ removes its identity and its singleton
stalk, contradicting the global bound through seventeen.
\end{proof}
